\section{Weighted square transport: ten positive variables of degree 190}
\label{sec:weighted-positive}

The next equation improves the degree of the inspected ten-positive-variable
family. Its exact profile is
\[
 (\text{variables},\text{degree},\text{binary operations})=(10,190,120).
\]
The operation split is 75 multiplications, 24 additions, and 21 subtractions;
squares and multiplications by constants are charged. This is a tradeoff
against the incoming centered free chart's profile $(10,206,90)$, and against
source178's scaled whole-product center $(10,208,116)$. No global minimum or
worldwide priority is asserted. The new step is the polynomial transport
weight $b^2UE$, which permits whole-product centering in the free chart
without excluding auxiliary Pell index three.

\subsection{The complete polynomial and its fixed data}

The independent variables are
\[
 J,F,\alpha,i,T,w,s,n,\eta,\zeta\in\mathbb Z_{>0}.
\]
The four fixed integers $(\beta,K_0,M_C,M_F)$ are the complete, unchanged
BFS943/radix67 tuple of the frozen source172 parent. The archive supplies all
four numerical payloads and reconstructs them from its complete finite
436-character recipe. The largest integer has 592,884 bits. This is the same
whole tuple throughout the arithmetic and logical calibration; its scalar
bounds alone do not specify the logical instance.

All letters on the following display are polynomial abbreviations, not
additional quantified variables:
\begin{align*}
 h&=\beta J,&q&=h+1,&C&=q-F-\alpha,\\
 X&=qw,&Y&=s,&\bar Y&=qY,\\
 k&=\eta+\zeta,&c&=kY+\eta,&A&=\bar Y(X+1)+1,\\
 \Delta&=A^2-1,&H&=4A-5,&v&=X\bar Y^2,\\
 \ell&=vk,&E&=1+\ell(\ell+k),&U&=ic^2,\\
 R&=[h(q-F)+\alpha](q^2-1)+(M_C+qM_F)J,&&&\\
 B&=(K_0+w)C-F,&a&=\Delta U^2E,&F_0&=1+a,\\
 G&=EF_0,&f&=1+2a,&Q&=4\Delta aF_0,\\
 b&=cT-R,&V&=fb-c,&K&=X(3c-2X)+2,\\
 e_6&=6(k-R-1)-\ell(k^2-4),&D&=6q^7,&M&=hDH.
\end{align*}
Define the following division-free center:
\begin{align*}
 e_2&=(k\bar Y)^2,&
 S&=(k\bar Y)\bigl[\ell+k-(k\bar Y)(\bar Y+1)\bigr],\\
 \theta&=U^2b,&
 R_c&=\theta(3E+e_2)-(b-c),\\
 S_c&=4E^2+3S^2+e_2^2,&
 \Omega&=U\bigl[2SR_c+A(4EV-\theta S_c)\bigr].
\end{align*}
Finally put
\begin{align}
 L&=b^2UE,\label{eq:weighted-L}\\
 Z_0&=hD\bigl[H(n-LB)+K\bigr]
        -H\bigl[DB+h(e_6-6q^4s)\bigr],\notag\\
 Z&=Z_0-M\Omega,\notag\\
 P&=Q\bigl(Z^2-16M^2GV^2\bigr)+16M^2G-Z^2.\label{eq:positive190}
\end{align}
The single equation is $P=0$. Every operation displayed is ordinary integer
addition, subtraction, or multiplication with fixed exponents. The literal
120-operation graph is supplied in \texttt{positive190/equation.json}.
There is no quotient, radical, variable exponent, factorial, sequence
predicate, or unknown coefficient inside this polynomial.

\paragraph{Consistency theorem.}
At the identified inherited whole-tuple calibration, ZFC proves
\[
 \bigl(\neg\exists J,F,\alpha,i,T,w,s,n,\eta,\zeta>0\;P=0\bigr)
 \quad\longleftrightarrow\quad\operatorname{Con}(\mathrm{ZFC}).
\]
Here consistency is finite syntactic non-refutability for the fixed effective
ZFC presentation of the parent. The result establishes an equivalence; it
does not establish either consistency or unsolvability separately.

\subsection{The inherited interface and its precise point of use}

The pinned source172 ordinary proof supplies the following named ingredients:
prime-valuation square extraction; adjacent-coefficient rational Pell descent;
signed all-index folding and normalized strong-index recovery; the full Lucas
expansion; the $\lambda$-uniform ratio and complete-product projection theorem;
and fresh native reconstruction followed by the whole-tuple syntactic
refutation calibration. The selected source text is retained as
\texttt{inherited/FREE\_CHART\_PROOF.md}, with its original source pins.
The new proof does not claim a freshly compiled transitive Lean/Rocq closure.

For proof notation only, let
\[
 t=\frac Bh,\quad \kappa=\frac KH,\quad
 \sigma=\frac{Y}{q^3},\quad e^*=\frac{e_6}{6},\quad
 \rho=\frac{e^*}{q^7},\quad z=\frac ZM,\quad W=1+hL.
\]
These are rational numbers until the appropriate divisibilities are proved.
Direct algebra gives
\begin{equation}
 z=n+\kappa+\sigma-Wt-\rho-\Omega.
 \label{eq:weighted-carrier}
\end{equation}
The elementary positive-domain bounds are
\[
 q\ge32,\quad X\ge q,\quad\bar Y\ge q,\quad A>q^2,
 \quad c\ge3,\quad U\ge c^2,\quad E>1,
 \quad a>c,\quad f>2c,\quad Q>c^2,\quad M>0.
\]
At a zero, use the exact identities
\[
 Q=4\Delta^2U^2G,\qquad \gcd(E,F_0)=1,
 \qquad (Q-1)Z^2=16M^2G(QV^2-1).
\]
The prime valuations show that $G=g^2$, $E=\tau^2$, and $F_0=f_0^2$ for
positive integers $g,\tau,f_0$; the added factor16 is a square. The retained
adjacent-coefficient descent applies to the initially rational
$y=|z|/(4g)$ and proves it is integral. Thus
\begin{equation}
 z=\varepsilon4gy,\quad\varepsilon\in\{-1,1\},\qquad
 QV^2-(Q-1)y^2=1.
 \label{eq:weighted-root}
\end{equation}
These arguments depend only on the final norm equation, not the carrier.
Signed folding and the first index bound then give, before any assumption on
the signs of $B,C,R,b,V$,
\begin{gather}
 c=\psi_A(p),\quad p\ge3\text{ odd},\qquad
 c\ge4\Delta+3>4q^4,\quad K>0,\notag\\
 b\ne0,\quad\operatorname{sgn}V=\operatorname{sgn}b,
 \quad y<\tfrac32|V|,\quad |V|<\tfrac32 f|b|,
 \quad |b|>4\Delta f-1.\label{eq:weighted-inherited-bounds}
\end{gather}
These bounds explicitly allow auxiliary index three. Since the final
quantities are integers, $|b|\ge4\Delta f$. Also
\[
 G=E+\Delta(UE)^2\le A^2(UE)^2,
 \qquad g\le AUE.
\]

\subsection{The center bound in the changed chart}

The whole-product center of source178 was developed in a different chart.
Its polynomial identity is valid here with its letter $Y$ replaced by
$\bar Y=qY$. The bound must be checked before recovering the sign of $B$.
Temporarily let $u$ be a formal indeterminate, and write
\begin{align*}
 E(u)&=e_2u^2+e_1u+e_0,\\
 e_1&=k^2\bar Y(1-2\bar Y-2\bar Y^2),\\
 e_0&=1+k^2\bar Y(\bar Y+1)
                   [\bar Y(\bar Y+1)-1],\\
 F(u)&=E(u)\bigl[(1+2(u^2-1)U^2E(u))b-c\bigr].
\end{align*}
For $j=1,3$, let $Q_j$ denote the ordinary polynomial quotient of $F(u)$ by
$u^j$, evaluated at $u=A$. This is formal polynomial division, never a floor
or rational instruction in the equation. Since
\[
 F(u)=(b-c)E(u)+2U^2b(u^2-1)E(u)^2,
\]
$Q_3$ is coefficientwise even. Exact polynomial algebra gives
\begin{equation}
 \Omega=U(4AEV-2Q_1-Q_3/2).
 \label{eq:weighted-quotient-center}
\end{equation}
The displayed division-free evaluator is precisely this polynomial. For
example, $S$ is the quotient of $E$ by $A$. If
\[
 J_1=S(2E-AS),\quad J_3=e_2(2S-Ae_2),\quad
 H_1=AE^2-J_1,\quad H_3=J_1-J_3,
\]
then $Q_1=(b-c)S+2\theta H_1$, $Q_3=2\theta H_3$, and
\[
 4H_1+H_3=A(4E^2+3S^2+e_2^2)-2S(3E+e_2).
\]
Substitution proves \eqref{eq:weighted-quotient-center} in the free integer
polynomial ring.

For a polynomial in $u$, use a superscript $\#$ for evaluation at $A$ after
taking absolute values of all coefficients. Put
$\delta=(\bar Y+1)/A<1/16$. The inequalities
\[
 |e_1|<2e_2(\bar Y+1),\qquad
 0<e_0\le e_2(\bar Y+1)^2
\]
give
\[
 E^\#(A)\le e_2(A+\bar Y+1)^2,
 \quad E\ge e_2A^2(1-2\delta),\quad E^\#(A)<2E.
\]
Consequently
\[
 F^\#(A)<2E(|b|+c)+8(A^2+1)U^2E^2|b|.
\]
At a zero $|b|\ge1$, $c+1<2U^2E$, and
$|V|>\Delta U^2E|b|$, because $f-c>a=\Delta U^2E$.
It follows that
\[
 \frac{F^\#(A)}{|EV|}
 <\frac4\Delta+\frac{8(A^2+1)}\Delta
 =8+\frac{20}\Delta<10.
\]
Deletion of low coefficients gives $|Q_j|\le F^\#(A)/A^j$. Thus, using
$A\ge16$ in \eqref{eq:weighted-quotient-center},
\[
 |\Omega-4AUEV|
 <10U|EV|\left(\frac2A+\frac1{2A^3}\right)<AU|EV|.
\]
In particular $\Omega$ has the sign of $V$, and the required uniform signed
estimate is
\begin{equation}
 3AUE|V|<|\Omega|<5AUE|V|.
 \label{eq:weighted-center-bound}
\end{equation}

\subsection{The new weighted transport argument}

The source178 weight $b^2$ does not have a proved adequate margin in this free
chart when auxiliary index three is possible. The additional factor $UE$ in
\eqref{eq:weighted-L} resolves that problem. Indeed, from
\eqref{eq:weighted-inherited-bounds} and
\eqref{eq:weighted-center-bound},
\begin{align}
 4gy+|\Omega|
 &<11AUE|V|
 <\frac{33}{2}AUEf|b|\notag\\
 &\le\frac{33A}{8\Delta}\,b^2UE
 <\frac L3.\label{eq:weighted-margin}
\end{align}
The last inequality follows from $99A<8(A^2-1)$ for $A\ge16$; here in fact
$A>1024$. Thus the proof neither assumes nor needs the stronger scaled-chart
index exclusion.

Since $k\ge2$, the inverse-Lucas correction satisfies
\[
 e^*\le k-R-1<b,
 \qquad \rho<\frac{|b|}{q^7}<\frac L2.
\]
This holds for every integer $R$, including zero and negative values.
If $B\le-1$, then $Wt=B/h+LB<-L$. Solving
\eqref{eq:weighted-carrier} for the positive input $n$ and using positive
$\kappa,\sigma$ gives
\[
 n=\varepsilon4gy+\Omega-\kappa-\sigma+Wt+\rho
 <\frac L3-L+\frac L2<0,
\]
a contradiction. Hence $B\ge0$. Since $F>0$ and $K_0+w>0$, the identity
$B=(K_0+w)C-F$ forces $C>0$.

Only now invoke the authenticated whole-tuple mask estimate
$3q<R<q^4-2$. Since $c>4q^4$, it implies $b,V>0$. The unchanged normalized
strong-index and signed auxiliary-index results identify $p=R$, from
\[
 V=-c\pmod f,\qquad V=-R\pmod c,\qquad c^2\mid2Ug.
\]
The retained parity argument yields odd $R$, even $q$, and odd $J$.

\subsection{Authentication, projection, and refutation}

The only changes in the rational carrier are the integer polynomial
$\Omega$ and the integer weight $L$, with $W\equiv1\pmod h$. The parent
authentication and projection proof therefore applies in the following order.
The first Pell norm gives $k=2\psi_{1+2v}(m)$, $m\ge2$, and monotonicity gives
$m<R$. Integral $z$ and the literal denominator give $q^4\mid e^*$.
The full retained Lucas identity is
\[
 3\bigl(e^*-(2m-R-1)\bigr)
   =4v^2(24Z_m-A_mS_m),
\]
where the right side is an integer expression from the finite binomial
expansion of the Pell sequence. Since $v=q^3ws^2$, at every prime dividing
$q$, including3, this implies $q^4\mid(2m-R-1)$. Its absolute value is below
$R<q^4$, so $m=(R+1)/2$.

Reduction of the literal numerator modulo $H$ gives $H\mid6hK$. The
$\lambda=q$ uniform ratio theorem and the complete-product strict bound force
$X=2^R$ and positive integral $\kappa$; this step does not cancel an
unproved factor of $6h$. Because $q\mid X$, the radix $q$ is a power of2.
Rank growth gives $q^6\mid X$, hence $q^8\mid v$. Only now does the same
Lucas identity imply $q^7\mid e^*$, so $\rho$ is integral.

Finally \eqref{eq:weighted-carrier}, with integral $\Omega$ and
$W\equiv1\pmod h$, gives $q^3\mid Y$ and $h\mid B$. Thus $t=B/h$ is a
nonnegative integer. These are exactly the parent native outer conditions.
The inherited fixed-minus parity, half-binomial rounding, and fresh native
reconstruction supply the native words at the complete fixed tuple; its
named syntactic calibration yields a finite ZFC refutation. This does not
replace that calibration by the elementary coefficient inequalities.

\subsection{Constructive positive completeness}

Conversely, a finite syntactic ZFC refutation supplies the parent canonical
outer data and positive $\eta,\zeta,i,T$. In particular $b,V>0$, $t\ge0$ is
integral, and $\rho,\kappa,\sigma$ are integral. The retained canonical
estimates give $g>\kappa+\sigma+|\rho|$. Keep these inputs and define
\[
 n=4gy+\Omega+(1+hL)t+\rho-\kappa-\sigma.
\]
Every term is integral; $\Omega>0$, $y\ge1$, and $t\ge0$ make $n>0$.
Its carrier satisfies $z=4gy$, so \eqref{eq:weighted-root} makes the literal
polynomial vanish. This proves both directions of positive solvability
versus finite refutability. The arithmetic arguments and the retained
proof-code equivalence are internalizable in ZFC; negation gives the
consistency theorem. No restriction to externally standard proof codes is
introduced.

\subsection{Exact degree and finite verification}

The free-chart degree weights for $(A,\bar Y,k,U,b,c)$ are
$(4,2,1,5,4,2)$. Exact sparse collection gives 84 center monomials, degree59,
and 594 monomials in
$D_{\rm norm}=\Omega^2-16GV^2$, of weighted degree 94. Its entire highest
component is
\[
 [D_{\rm norm}]_{94}
 =8A^6\bar Y^8k^8U^6b^2(A^2-6\bar Y^4).
\]
The exact output decomposition is
\[
 P=Q\bigl[Z_0^2-2MZ_0\Omega+M^2D_{\rm norm}\bigr]
      +16M^2G-(Z_0-M\Omega)^2.
\]
Here
\[
 \deg(Q,M,G,V,E,\Omega,L,Z_0)=(72,12,46,36,14,59,27,41).
\]
The five respective degree bounds in the decomposition are
$154,184,190,70,142$, so the norm-defect term uniquely controls the top
degree. With $q_0=\beta J$ and $k_0=\eta+\zeta$, the complete actual
highest component is
\begin{equation}
 [P]_{190}
 =9\cdot2^{11}q_0^{74}w^{18}s^{52}i^{10}k_0^{32}
       (q_0-F)^2(w^2-6s^2).
 \label{eq:positive190-leading}
\end{equation}
This is nonzero for the actual fixed positive $\beta$, proving exact degree 190.
The full leading form has 198 collected monomials.

The independent standard-library checker verifies the actual literal graph,
all-variable/port/gate liveness, the entire center and norm-defect identities,
the complete leading form, all four payload hashes, and 60 deliberately
corrupted center operators. Both normal and optimized Python runs pass.
The separate frozen numeral recipe reconstructs every coefficient byte.
These checks authenticate finite artifacts; the unbounded Pell and logical
theorems remain ordinary mathematical proofs with the stated inherited
dependencies.

\paragraph{Further research.}
The weighted transport argument admits other integer polynomial centers with
the same uniform $O(AUE|V|)$ signed bound. Reducing the 30-gate center
evaluator is a concrete cost target. Reducing the degree requires more than
the next term of the same $A$-Laurent expansion: that term removes the
$8A^8\bar Y^8k^8U^6b^2$ component but retains a nonzero multiple of
$A^6\bar Y^{12}k^8U^6b^2$, of the same weighted degree 94. Any proposed
division by $A$ or $w$ must be replaced by a proved polynomial identity or
an explicitly counted additional witness.
